% Appendix A. Grounded rewrite 2026-11-03 (branch docs/cep-grounded-draft): lists the scripts that actually generate the
% paper's numbers and figures. Test counts from `uv run pytest --collect-only`.
\section{Reproduction}\label{app:repro}

Every input value lives in \code{configs/camera_100.toml}, tagged with its source class; \code{tests/test_provenance.py}
fails if a value lacks a class. The environment and the optics tests are set up with
{\footnotesize
\begin{verbatim}
uv sync --group dev --group optics \
  --group annotator
uv run pytest -q tests/test_optics.py \
  tests/test_scale_chain.py \
  tests/test_provenance.py   # 30 tests
\end{verbatim}}
\noindent The full suite collects 142 tests (\code{uv run pytest -q}). \tabref{tab:repro} lists the generators.
\code{scripts/rerun_paper_numbers.py} writes each computed number of \S\S\,\ref{sec:optics}--\ref{sec:csd} with its
previous value, source class and generating function to \code{docs/optics/rerun_results.md};
\code{scripts/detector_checks.py} writes the label-free detector checks to \code{docs/optics/detector_checks.md};
\code{optics.scale_chain} writes the scale budget to \code{docs/optics/scale_chain.md}; \code{make -C papers/cep
figures} collects the figures. The recordings are read from \code{datasets/}.

\begin{table}[!ht]
\centering\scriptsize\setlength{\tabcolsep}{3pt}
\caption{Generators of the paper's numbers and figures. Each command follows \texttt{uv run python}
(\texttt{uv run --frozen python} for the rerun script, \texttt{uv run --with openpyxl python} for the anchors);
$^\ast$: runs without the videos.}\label{tab:repro}
\begin{tabular}{@{}>{\raggedright\arraybackslash}p{0.5\columnwidth}>{\raggedright\arraybackslash}p{0.46\columnwidth}@{}}
\toprule
Command & Output \\
\midrule
\code{-m optics.scale_chain} (\code{--no-video}$^\ast$: budget) & \secref{sec:chain}, \tabref{tab:scale} \\
\code{-m optics.error_budget}$^\ast$ & 1$\sigma$ budget, \secref{sec:slab} \\
\code{scripts/rerun_paper_numbers.py} & other numbers of \S\S\,\ref{sec:optics}--\ref{sec:csd};
Tables~\ref{tab:dof}, \ref{tab:crop}, \ref{tab:ws-compare} \\
\quad\code{--no-video}$^\ast$ / \code{--fig07} & optics rows / \figref{fig:overlays} \\
\code{scripts/check_clip_overlap.py} & clip date, \S\S\,\ref{sec:chain}--\ref{sec:recording} \\
\code{scripts/check_anchors.py} & anchors, clock, \secref{sec:expected} \\
\code{scripts/microscope_summary.py}$^\ast$ & \tabref{tab:microscope} \\
\code{scripts/detector_checks.py} & \tabref{tab:ws-sweep}, \figref{fig:time} \\
\code{-m optics.plots}$^\ast$ & Figs.~\ref{fig:coc},
\ref{fig:sharpness}, \ref{fig:dof} \\
\code{-m optics.optical_layout}$^\ast$ & \figref{fig:layout} \\
\code{-m optics.measure} & \figref{fig:measured} \\
\code{-m optics.qualitative} & \figref{fig:detection} \\
\code{-m optics.disambiguate} & Figs.~\ref{fig:discrim}, \ref{fig:disamb} \\
\code{-m optics.csd} & \figref{fig:csd} \\
\bottomrule
\end{tabular}
\end{table}

\figref{fig:frames} is the first frame of the case-study clip as stored in the repository; it has no generator
script yet.
